\section{The verified 215-state realization of the revised scheduler}
\label{sec:scheduler-realization}

The revised scheduler has a complete 215-state binary realization in
\texttt{q13\_cb215.native.tm}. Its source is
\texttt{q13\_cb.sites.json}. This is a calibrated intermediate construction;
its count does not meet the requested threshold of 214 states.

The exact artifact hashes are
\begin{center}\ttfamily\footnotesize
source: 79425cf76b24b8739fd00e3fd091ca00de3b3379324f941158ba5e58061dc2b0\\
table: c6ff22b5446bb7f14a9f8ea02828fa694291ce1b7c71cf1ea3cb488d03d30dc5.
\end{center}
The source contains 4,752 primitive occurrences and 202 distinct source
DAG nodes. Its maximum address depth is 47 and maximum zero run is six.
The runtime receipt proves all 6,177 guarded primitive cases, including
arbitrary natural primary values, arbitrary natural gap, and an arbitrary
nonempty finite spare list, with complete zero exterior. It reports
215 working states and no failed cases. This is an ordinary mathematical
runtime certificate, not a proof-assistant kernel artifact.

Blank execution reaches dispatcher state $84$ after $88{,}116$ binary
transitions, with head at $-12$, gap $5$, 45 unit-valued spare words, and
logical register tuple
\[
 (N,T,W,A,D,H,X,Y,L)=(5,0,0,0,0,0,0,0,4).
\]
The complete tape is specified by the startup receipt, including its
zero exterior. The nonzero list does not require a new soundness
assumption: every theorem register and decoded stack head is zero.

There is a particularly direct reduction of this actual seed to the
accepted K215 seed. The first $G$ decrement changes $L=4$ to $3$.
Since $P(2,0)=3$, the next decoder selects $\mathrm{MP}$, and its
remaining field is zero. The command is therefore $\mathrm{MP}(0)$.
Its comparison succeeds because $I(0,0)=0$ is the cached value, but its
new top is zero and its new stack is $P(0,0)=0$; every other retained
register stays zero. It returns after 65 source primitives to
\[
 (5,0,0,0,0,0,0,0,0).
\]
The portable source audit verifies this exact finite reduction. The
counter-offset relation in Section~\ref{sec:refill-scheduler} now relates
this root directly to the accepted source's root
$(4,0,0,0,0,0,0,0,0)$. All subsequent command execution is matched by
the proved stuttering simulation.

Consequently the whole-tape runtime certificate, finite blank prefix,
source simulation, and inherited calculus calibration give the ordinary
ZFC theorem
\[
 \neg\operatorname{Halt}(\mathrm{q13\_cb215},\mathrm{blank})
 \quad\longleftrightarrow\quad
 \operatorname{Con}(\mathrm{ZFC}).
\]
The construction exposes a different state-allocation opportunity:
the source no longer needs state 13 as its unique decoder for the refill
continuation. State 13 is still needed by the present engine realization,
so deleting it requires an additional valid engine replacement.
